\begin{definition}[Strong Convergence]
    A sequence \( ({x}_{n}) \) in a normed space \( X  \) is said to be \textbf{strongly convergent} (or \textbf{convergent in the norm}) if there is an \( x \in X  \) such that 
    \[  \lim_{ n \to \infty  }  \|{x}_{n} - x \| = 0.  \]
\end{definition}

\begin{definition}[Weak Convergence]
    A sequence \( ({x}_{n}) \) in a normed space \( X  \) is said to be \textbf{weakly convergent} if there is an \( x \in X  \) such that for every \( f \in X' \),
    \[  \lim_{ n \to \infty  } f({x}_{n}) = f(x). \]
    This is written 
    \[  {x}_{n} \xrightarrow{w} x.   \] 
    The element \( x  \) is called the \textbf{weak limit} of \( ({x}_{n}) \), and we say that \( ({x}_{n}) \) \textbf{converges weakly to} \( x \).
\end{definition}

\begin{lemma}[Weak Convergence]
    Let \( ({x}_{n}) \) be a weakly convergent sequence in a normed space \( X  \), say, \( {x}_{n} \xrightarrow{w} x  \). Then: 
    \begin{enumerate}
        \item[(a)] The weak limit \( x  \) of \( ({x}_{n}) \) is unique;
        \item[(b)] Every subsequence of \( ({x}_{n}) \) converges weakly to \( x  \);
        \item[(c)] The sequence \( (\|{x}_{n}\|) \) is bounded.
    \end{enumerate}
\end{lemma} 

\begin{theorem}[Strong and Weak Convergence]
    Let \( ({x}_{n}) \) be a sequence in a normed space \( X  \). Then:
    \begin{enumerate}
        \item[(a)] Strong convergence implies weak convergence with the same limit;
        \item[(b)] The converse of (a) is not generally true;
        \item[(c)] If \( \text{dim}(X) < \infty  \), then weak convergence implies strong convergence.
    \end{enumerate}
\end{theorem}

\begin{lemma}[Weak Convergence]
    In a normed space \( X  \), we have \( {x}_{n} \xrightarrow{w} x  \) if and only if  
    \begin{enumerate}
        \item[(A)] The sequence \( (\|{x}_{n}\|)  \) is bounded;
        \item[(B)] For every element \( f  \) of a total subset \( M \subseteq  X' \), we have \( f({x}_{n}) \to f(x) \).
    \end{enumerate}
\end{lemma}
